% Options for packages loaded elsewhere
\PassOptionsToPackage{unicode}{hyperref}
\PassOptionsToPackage{hyphens}{url}
%
\documentclass[
]{article}
\usepackage{amsmath,amssymb}
\usepackage{iftex}
\ifPDFTeX
  \usepackage[T1]{fontenc}
  \usepackage[utf8]{inputenc}
  \usepackage{textcomp} % provide euro and other symbols
\else % if luatex or xetex
  \usepackage{unicode-math} % this also loads fontspec
  \defaultfontfeatures{Scale=MatchLowercase}
  \defaultfontfeatures[\rmfamily]{Ligatures=TeX,Scale=1}
\fi
\usepackage{lmodern}
\ifPDFTeX\else
  % xetex/luatex font selection
\fi
% Use upquote if available, for straight quotes in verbatim environments
\IfFileExists{upquote.sty}{\usepackage{upquote}}{}
\IfFileExists{microtype.sty}{% use microtype if available
  \usepackage[]{microtype}
  \UseMicrotypeSet[protrusion]{basicmath} % disable protrusion for tt fonts
}{}
\makeatletter
\@ifundefined{KOMAClassName}{% if non-KOMA class
  \IfFileExists{parskip.sty}{%
    \usepackage{parskip}
  }{% else
    \setlength{\parindent}{0pt}
    \setlength{\parskip}{6pt plus 2pt minus 1pt}}
}{% if KOMA class
  \KOMAoptions{parskip=half}}
\makeatother
\usepackage{xcolor}
\setlength{\emergencystretch}{3em} % prevent overfull lines
\providecommand{\tightlist}{%
  \setlength{\itemsep}{0pt}\setlength{\parskip}{0pt}}
\setcounter{secnumdepth}{-\maxdimen} % remove section numbering
\ifLuaTeX
  \usepackage{selnolig}  % disable illegal ligatures
\fi
\IfFileExists{bookmark.sty}{\usepackage{bookmark}}{\usepackage{hyperref}}
\IfFileExists{xurl.sty}{\usepackage{xurl}}{} % add URL line breaks if available
\urlstyle{same}
\hypersetup{
  hidelinks,
  pdfcreator={LaTeX via pandoc}}

\author{}
\date{}

\begin{document}

Neural Representations and Benign-Only Anomaly Detection for Unknown
Attacks in In-Vehicle CAN Traffic Adda Surendra, Gera Vignan, and
Vundadi Mani Kumar Abstract---An in-vehicle intrusion detector is often
trained with a finite set of attack labels, whereas a deployed vehicle
may encounter traffic patterns that were not represented during
training. This study examines one practical response to that mismatch:
train a multilayer perceptron (MLP) to classify known Controller Area
Network (CAN) traffic classes, use its 64- dimensional hidden
activations as a representation, and fit an Isolation Forest (IF) to
benign training representations. The experiments use the public
CAN-train-and-test v1.5 data and treat six attack families in an
evaluation capture as unknown to the six-class MLP. At the primary
operating point, defined by the 99th percentile of benign calibration
scores, the detector obtains 30.68\% precision, 19.00\% recall, and a
23.47\% F1-score, with a 1.04\% false-positive rate. Its continuous
scores yield ROC- AUC 0.9423 and PR-AUC 0.4299. The family-wise results
are uneven: 65.27\% of FUZZING records are detected, while no DOUBLE,
INTERVAL, SPEED, or TRIPLE records cross the primary threshold. A
97.5th-percentile operating point produces 49.37\% F1 and 62.53\%
recall, but it is retained as exploratory because the same unknown-test
file had been examined during earlier threshold analysis. The comparison
with an Isolation Forest fitted to the original features shows a
distinction between ranking and thresholded decisions: the learned
representation improves ROC-AUC and PR-AUC, but not F1 at the primary
threshold. These results do not establish a new general-purpose detector
or a complete reproduction of FedGAN-ID. They provide a reproducible
account of one MLP--IF pipeline, its operating- point sensitivity, and
the attack families it fails to identify. Index Terms---Automotive
cybersecurity, Controller Area Net- work, in-vehicle network, intrusion
detection, unknown attacks, anomaly detection, neural representation,
Isolation Forest. I. INTRODUCTION A modern passenger vehicle is not a
single computer. It is a collection of electronic control units (ECUs)
that exchange short messages to coordinate functions such as engine
management, braking, steering, body electronics, and driver assistance.
The Controller Area Network (CAN) is one of the communication
technologies used for this exchange. Its design favours timely, compact
broadcast communication rather than the security properties normally
associated with networks exposed to untrusted participants. In
particular, a classical CAN frame does not by itself authenticate the
ECU that transmitted it. An attacker who gains a route onto the The
authors are with the Department of Computer Science and En- gineering
(Data Science and Artificial Intelligence), Indian Institute of
Information Technology Design and Manufacturing, Kancheepuram, India
(e-mail: surendraadda@gmail.com; gvignankv@gmail.com; vundadimaniku-
mar@gmail.com). This project was carried out under the guidance of Dr.
Bhale Pradeepkumar Gajendra, Assistant Professor, Department of Computer
Science and Engineering, IIITDM Kancheepuram. in-vehicle bus may
therefore attempt to inject frames, replay observed messages, or
impersonate the traffic of another component. The security implications
of such access have been demonstrated in automotive security research
and motivate monitoring at the network level {[}1{]}--{[}3{]}. An
intrusion detection system (IDS) can complement pre- ventive controls by
observing bus activity and reporting be- haviour that departs from an
expected profile. CAN IDS research has considered several kinds of
evidence. Some methods examine message identifiers and payload bytes;
oth- ers model frame frequency, inter-arrival time, message or- dering,
or ECU-specific timing characteristics. These signals are related, but
they are not interchangeable. A denial-of- service attack may create a
pronounced rate change, whereas a masquerade or replay-like pattern may
preserve a familiar identifier or approximate a legitimate transmission
schedule. Timing-based methods can be useful when normal periodicity is
stable, while content-oriented approaches can detect payload changes
that a rate-only monitor would not see {[}2{]}--{[}4{]}. A common
experimental assumption is that the attack classes available during
training are also the classes that matter at evaluation time. That
assumption is convenient for closed-set classification, but it does not
describe the full operational problem. Attack variants change, and a
vehicle manufacturer may not possess labelled examples of every
malicious behaviour before deployment. An IDS intended to recognise
previously unseen attacks must therefore do more than assign a test
message to one of the labels learned during training. It needs a
mechanism for identifying observations that are sufficiently unlike the
benign operating region, while keeping false alarms at a usable level.
This paper studies a small, testable pipeline for that setting. A
supervised MLP is trained on six known classes. The penultimate hidden
activation is then treated as a compact representation of each CAN
record. An Isolation Forest is fitted only to the hidden representations
of benign training records. A threshold is estimated from benign
calibration records and subsequently applied to an evaluation file
contain- ing benign traffic and six attack families excluded from MLP
training. The design is deliberately narrower than the system that
motivated the work. FedGAN-ID combines CAN-message graph construction, a
GAN-based detector, federated learning, differential privacy, and
energy-based unknown-attack detec- tion {[}5{]}. The present
implementation does not include those graph, GAN, federated, or privacy
mechanisms; it examines an MLP-representation/IF alternative on a public
The central question is not whether the pipeline can be made to report a
high aggregate score at any threshold. It is whether a representation
trained for known-class discrimination retains useful structure for an
anomaly detector, and whether this structure is useful across different
unseen attack families. Four questions guide the experiments: 1) Does
the MLP hidden representation provide a useful anomaly ranking when the
IF is fitted using benign records only? 2) How do precision, recall,
F1-score, and false-positive rate change between the initially selected
99th-percentile threshold and a more permissive exploratory threshold?
3) Do the six unknown attack families exhibit similar detection
behaviour, or is the aggregate result dominated by a subset of them? 4)
How does the representation-based detector compare with an IF fitted
directly to the 13 input features and with the executed energy-based OOD
experiments? The contributions are empirical rather than architectural.
First, the manuscript documents the data preparation and unknown-attack
protocol used in the project. Second, it reports the complete primary
confusion matrix and per-family recalls, not accuracy alone. Third, it
separates the preselected operat- ing point from the exploratory
threshold that was examined after test-set inspection. Fourth, it
analyses why a strong ROC-AUC can coexist with a weak F1-score at a
fixed threshold. Finally, it records negative results from the energy-
based experiments and makes explicit which components of FedGAN-ID were
not implemented. This narrower framing is important: the results should
be judged as an evaluation of the implemented pipeline, not as a claim
to reproduce or outperform a substantially different system. II.
TECHNICAL BACKGROUND AND RELATED WORK A. CAN Communication and Its
Security Context CAN was designed for communication among embedded
controllers under constraints on message size, bus access, and latency.
The arbitration mechanism prioritises messages by identifier, and frames
are broadcast to nodes on the bus. These design choices support robust
real-time communication, but the protocol's basic frame format does not
provide crypto- graphic source authentication. Consequently, a receiving
ECU generally cannot establish the sender's identity solely from the CAN
identifier. A compromised node or an external interface that reaches the
bus can potentially transmit a frame using an identifier associated with
a legitimate signal. The Bosch CAN specification describes the
protocol's communication format and bus access model; it should not be
confused with a security mechanism that authenticates each sender
{[}1{]}. The absence of native source authentication does not imply that
every CAN message is equally easy to forge. Real vehicle traffic
contains regularities: identifiers tend to be associated with particular
message roles, payload fields often evolve within constrained ranges,
and many messages are transmitted periodically. An IDS can learn a
baseline from these regu- larities. The difficulty is that the same
regularities can vary with vehicle state, driving conditions, diagnostic
activity, and vehicle model. A detector that memorises a single capture
can raise alarms on legitimate but unfamiliar operation. This makes the
choice of data split and threshold policy as important as the classifier
architecture. B. Families of CAN Intrusion Detectors A useful way to
organise CAN IDS designs is by the evidence they analyse.
Content-oriented systems examine the CAN identifier, payload length, or
data bytes. Statistical sys- tems summarise the frequency or
distribution of identifiers and payload values. Temporal systems monitor
inter-arrival time or the periodicity of messages. Sequence-oriented
systems model transitions across multiple frames. Physical-layer and
ECU- fingerprinting approaches use characteristics such as clock
behaviour or signal properties to associate messages with particular
transmitters. These approaches address related but distinct questions:
whether a frame looks unusual, whether its timing is unusual, and
whether it came from an expected ECU. Surveys of AI-based in-vehicle
IDSs discuss the range of learning strategies and the difficulty of
comparing results across datasets and evaluation protocols {[}2{]},
{[}6{]}. Muter and Asaj investigated entropy-based anomaly de- tection
for in-vehicle networks, illustrating the use of traffic statistics as a
baseline for abnormality {[}4{]}. Cho and Shin's clock-based IDS uses
periodic message timing to fingerprint ECUs and detect deviations in
clock behaviour {[}3{]}. Such work motivates the inclusion of timing
features, but it also illustrates a limitation of record-wise models:
when a detector processes each row independently, information encoded in
longer mes- sage sequences or transmitter-specific clock patterns may be
only partially represented. C. Known-Attack Classification and
Unknown-Attack Detec- tion A closed-set classifier estimates a label
from a fixed class list. If the input belongs to a class absent from
training, the classifier may still return one of its known labels with
high confidence. This is not necessarily a software defect; it follows
from the classification objective, which selects among the output
classes rather than explicitly modelling every possible unknown class.
Unknown-attack detection therefore requires an additional rule, score,
or model that can indicate when an observation falls outside the learned
benign region. Several strategies are possible. A one-class model can be
trained on benign examples alone. A classifier can use confi- dence or
energy scores as an indicator of out-of-distribution input. A generative
model can approximate normal traffic and use a reconstruction or
likelihood-related score. A hybrid detector can combine several types of
evidence. FedGAN- ID is an example of a more elaborate design: its
published description uses CAN-message graphs, a GAN-based detec- tor, a
second-level attack classifier, federated learning across vehicles,
differential privacy, and energy-based OOD detection {[}5{]}. It
evaluates its method on three real vehicles supported by XPeng. These
system-level properties are not assumed in this work. Page 4 of 14 - AI
PROJECT MANUSCRIPT, NOVEMBER 2026 3 D. Why Use a Neural Representation
Before Isolation Forest? Isolation Forest isolates samples through
random partition- ing. Anomalous observations are expected, on average,
to require fewer splits to isolate than observations in dense regions.
The original method does not require attack labels during fitting and is
therefore a plausible choice when only benign examples are used to
define the reference region {[}7{]}. However, Isolation Forest applied
to raw features can be sensitive to mixed scales and to the geometry of
the original feature space. A neural network trained for known-class
classification transforms input features into a hidden space shaped by
the classification objective. In principle, that space may bring
together records that share known-class structure and separate records
whose combinations of features differ. The assumption is not guaranteed
to hold for unseen attacks. A representation optimised to discriminate
six known labels may discard varia- tion that does not help distinguish
those labels. It can also map an unseen attack close to a familiar
benign or attack cluster. The MLP--IF combination is thus an empirical
question: the representation may improve the ordering of benign and
attack scores, but threshold performance and family coverage must be
measured separately. E. Dataset Generalisation as an Evaluation Concern
The can-train-and-test dataset was curated to support eval- uation
across attack types and vehicles, and its associated publication
discusses the limitations of available labelled auto- motive datasets
{[}8{]}. The existence of multiple capture subsets is valuable because a
randomly shuffled row-level split may place nearly identical traffic
conditions on both sides of an evaluation. At the same time, a
particular test capture does not automatically establish generalisation
to all vehicle platforms or all driving conditions. The present study
uses the named unknown-attack evaluation file and makes no claim that
every aspect of cross-vehicle generalisation has been tested. III. STUDY
DESIGN AND PROBLEM FORMULATION A. Threat Model The detector is intended
to identify CAN records whose extracted attributes are inconsistent with
the learned benign region. The evaluation assumes that the attacker can
cause one or more non-benign CAN records to appear in the observed
stream, and that the monitoring component can access the identifier,
payload length, payload bytes, and capture times- tamps required to
derive the input features. The experiment does not model cryptographic
key compromise, ECU firmware compromise, physical-layer signal
manipulation, or an adver- sary that deliberately optimises messages
against the detector. It also does not establish which ECU transmitted a
suspicious frame. The output is a binary record-level anomaly flag, not
an attribution or attack-family diagnosis. The threat model is
intentionally limited to the capabilities represented by the dataset
labels. The six evaluation families are treated as unknown because they
are not among the six MLP training labels. This is an operational
definition of unknown relative to the training protocol, not a claim
that the attack techniques themselves are novel in the broader security
literature. B. Notation and Learning Objective Let each captured record
be represented by a feature vector xi ? R13 and a known training label
yi ? \{1, . . . , 6\}. The six labels are BENIGN, DoS, FORCE NEUTRAL,
FUZZING TARGETED, RPM, and STANDSTILL. The MLP learns a parameterised
mapping g? : R13 -> R6 for this closed set. We denote the penultimate
64-dimensional activation by zi = f?(xi), zi ? R64 . (1) The IF is
fitted to \{zi : yi = BENIGN\} only. It produces a scalar score s(z)
whose sign is oriented in the implementation so that larger values
indicate a more anomalous record. The decision rule is \textasciicircum{} a(x) = ?
{[}s(f?(x)) \textgreater{} ?{]} , (2) where \textasciicircum{} a = 1 denotes an anomaly
and ? is estimated from benign calibration scores. The evaluation target
is binary: BENIGN records are nega- tive, and records from any of the
six held-out attack families are positive. This aggregation is useful
for measuring whether a previously unseen attack record is flagged, but
it does not mean that all attacks are behaviourally equivalent. The per-
family analysis is therefore reported alongside the aggregate confusion
matrix. C. Research Questions and Interpretation Boundary The study is
designed to answer the four research questions introduced above, under
the following interpretation boundary. A higher ROC-AUC indicates that
positive records tend to receive higher scores than negative records
across score thresh- olds. It does not specify which threshold should be
deployed. F1-score, precision, recall, and false-positive rate depend on
the selected threshold and the mix of records in the evaluation set. The
main result is consequently tied to a preselected calibration rule, and
exploratory alternatives are kept separate. This distinction is
especially relevant in an imbalanced test set, where a high proportion
of benign records can make accuracy appear favourable even when a
substantial fraction of attacks is missed. Precision--recall analysis is
valuable in such settings because it makes the cost of false alerts and
missed positives more visible {[}9{]}. IV. DATASET AND DATA PREPARATION
A. Source Dataset and Capture Selection The experiments use
CAN-train-and-test v1.5, an openly released dataset intended for
developing and assessing auto- motive IDSs {[}8{]}, {[}10{]}. The
repository is maintained by the dataset authors and provides multiple
train/test subsets. The dataset paper describes a curation strategy
intended to test generalisation across attack types and vehicles; this
project uses only one configured subset of that broader resource. The
project uses the set 01 training data to construct the Page 5 of 14 - AI
PROJECT MANUSCRIPT, NOVEMBER 2026 4 known-class training and calibration
files, and the designated unknown-attack test capture for evaluation. In
the local project workflow, the unknown evaluation file is referred to
as test\_ 03\_known\_vehicle\_unknown\_attack. The filename's
distinction is retained: the experiment treats its attack families as
unknown to the classifier; it does not label the evaluation as an
unseen-vehicle test. Three processed files are used. The training file
is used to fit the scaler and MLP, and its benign subset is used to fit
the IF. The calibration file is used to estimate the anomaly- score
threshold from benign records only. The unknown-test file is used to
calculate the final metrics and family-wise recall. The labels from the
unknown-test file are used for evaluation, not as model-fitting targets.
However, the evaluation file was examined during earlier exploratory
threshold analysis. For this reason, the 97.5th-percentile result is not
presented as an independent confirmation. B. Record Counts and Class
Composition Table I summarises the processed records. The training file
contains 461,345 rows, of which 433,538 are benign and 27,807 are known
attacks. The calibration file contains 492,435 rows; 463,134 of these
are benign. It also contains known attack records, but the threshold
calculation uses only its benign subset. The unknown-test file contains
882,269 rows, including 861,444 benign records and 20,825 records from
the six attack families. The six known attack counts sum to 27,807 in
the training file. In the unknown-test file, the attack count is 20,825,
or approximately 2.36\% of all processed records. This prevalence is a
property of the prepared evaluation file, not an estimate of the
frequency of attacks in normal vehicle operation. During preprocessing,
benign records were subsampled at 10\% per capture, subject to a maximum
of 100,000 retained benign records per file, while attack records were
retained. This changes the class mixture and therefore affects metrics
such as precision and accuracy. The counts in the table should be
retained when results are compared or reproduced. The calibration
partition includes records from known attack classes, but only records
whose label is BENIGN are passed to the threshold calculation in the
embedding-IF script. The benign-only selection is relevant: using
attack-labelled calibra- tion records to choose a threshold would change
the evaluation protocol and would no longer represent a benign-reference
calibration procedure. C. Feature Construction Each row is described by
13 numeric features. The CAN identifier and payload length provide
frame-level information. The eight payload byte fields preserve the byte
values as separate numeric inputs. Three timing features capture local
temporal context: the inter-arrival time between the current record and
its predecessor, the inter-arrival time since the previous message with
the same identifier, and a rolling mean of inter-arrival time. Absolute
timestamps are not passed to the MLP or the IF. Temporal features are
calculated in capture order within the original capture file before the
processed partitions are consumed by the model. Negative time
differences are clipped to zero in preprocessing. This preserves the
chronological ordering available in the capture, but it does not
construct an explicit multi-frame sequence tensor. Once calculated, each
row is treated as one 13-dimensional observation. Thus, the model can
use local timing summaries but cannot directly learn arbitrary
long-range message dependencies. D. Feature Scaling and Label Mapping A
standard scaler is fitted to the training feature matrix. The same
fitted transformation is applied to calibration and test rows. No
calibration or test feature values are used to estimate the scaler
parameters. The MLP label mapping is restricted to the six known
training classes. For binary unknown-attack evaluation, the ground-truth
label is converted to zero for BENIGN and one for every non-BENIGN
record. The six unknown labels are never added to the MLP output layer.
E. Preprocessing Risks Preprocessing itself can affect apparent model
performance. Subsampling benign records changes the empirical class
prior. Calculating timing features within each original capture avoids
using a timestamp from one capture as the predecessor of a record from
another capture; however, the exact capture boundaries and ordering must
be preserved to reproduce the features. Numeric treatment of CAN
identifiers also imposes a geometry on the identifier values. The
numerical distance between two IDs is not necessarily a meaningful
semantic distance between their message functions. The current pipeline
uses this numeric encoding because it is what the implemented model
consumes; it does not claim that CAN identifiers are naturally
continuous variables. V. PROPOSED EXPERIMENTAL PIPELINE A. Pipeline
Overview Fig. 1 shows the stages actually implemented. Raw capture files
are first converted into processed rows. The scaler and six-class MLP
are fitted using the training partition. The penultimate activation is
extracted for benign training rows, benign calibration rows, and all
unknown-test rows. The IF is fitted to the benign training activations.
Its anomaly scores on benign calibration activations determine the
primary threshold. The decision rule is then applied to the unknown-test
scores. The data flow contains two different learning objectives. The
MLP is supervised: its purpose is to discriminate among the known
labels. The IF is unsupervised with respect to attacks: it is fitted
only to benign representations and does not use known-attack labels. The
two objectives should not be conflated. In particular, the IF is not
trained to classify DOUBLE, FUZZING, INTERVAL, SPEED, SYSTEMATIC, or
TABLE I PROCESSED DATASET COMPOSITION USED IN THE EXPERIMENTS Partition
Label/group Records Role in experiment Training BENIGN 433,538 MLP
training and IF fitting Training DoS 9,139 MLP known class Training
FORCE NEUTRAL 6,500 MLP known class Training FUZZING TARGETED 6,615 MLP
known class Training RPM 3,338 MLP known class Training STANDSTILL 2,215
MLP known class Training Total 461,345 Calibration BENIGN 463,134
Threshold estimation Calibration Known attack records 29,301 Not used
for threshold Calibration Total 492,435 Unknown test BENIGN 861,444
Negative evaluation class Unknown test DOUBLE 1,292 Unknown attack
family Unknown test FUZZING 5,975 Unknown attack family Unknown test
INTERVAL 4,418 Unknown attack family Unknown test SPEED 5,321 Unknown
attack family Unknown test SYSTEMATIC 1,043 Unknown attack family
Unknown test TRIPLE 2,776 Unknown attack family Unknown test Total
882,269 TABLE II INPUT FEATURES Feature Meaning in the processed row
can\_id Arbitration identifier represented numerically payload\_length
Number of data bytes in the frame byte\_0--byte\_7 Eight payload byte
values; unused byte positions follow the preprocessing convention
delta\_t\_ms Time since the preceding record in capture order
same\_id\_delta\_t\_ms Time since the previous occurrence of the same
CAN identifier rolling\_mean\_delta\_t\_ms Rolling mean of the
inter-arrival-time feature Fig. 1. Implemented MLP-representation and
benign-only Isolation Forest pipeline. The MLP uses the six known
labels; the IF fitting and primary threshold estimation use benign
records only. B. Multilayer Perceptron The MLP receives the 13
standardised features. Its fully connected network has a first hidden
layer with 128 units, a ReLU activation, and dropout of 0.2; a second
hidden layer with 64 units, another ReLU activation and dropout of 0.2;
and a six-unit linear output layer. The architecture can be written as
h1 = Dropout0.2(ReLU(W1x + b1)), (3) h2 = Dropout0.2(ReLU(W2h1 + b2)),
(4) l = W3h2 + b3, (5) where l ? R6 contains the class logits. During
embedding extraction the network is placed in evaluation mode, so
dropout is disabled, and the 64-dimensional activation immediately after
the second ReLU is captured. In notation, this activation is z = f?(x).
The training script uses cross-entropy loss with square- root
inverse-frequency class weights. If nc is the number of training records
for class c and N is the total number of training records, the
unnormalised weight is proportional to ? N/(6nc); the weights are
normalised to have mean one. This reduces the contribution imbalance
caused by class sizes without applying the full inverse-frequency
weighting. Optimisation uses Adam with learning rate 0.001, batch size
4096, and 25 epochs. The random seed is set to 42 for the Python
numerical library and the PyTorch model. The MLP Page 7 of 14 - AI
PROJECT MANUSCRIPT, NOVEMBER 2026 6 is run on CPU in the recorded
experiment. The training log reported a final training accuracy of
99.86\%. This figure is a training-set statistic. It does not measure
unknown-attack detection performance and must not be presented as a
held-out classification accuracy. The down- stream anomaly detector is
evaluated separately using the unknown-test file. C. Extraction of the
Hidden Representation The activation after the second ReLU is a
64-dimensional vector. The dropout layer follows this activation during
MLP training, but embedding extraction is performed with the model in
evaluation mode. Consequently, the IF receives the deterministic hidden
activation from the trained model rather than the output logits. This
choice makes the representation a learned feature map from the original
13-dimensional space to a 64-dimensional space. The representation has
no explicit metric-learning loss or contrastive objective. Its geometry
is induced by the six- class cross-entropy objective. This is a design
choice, not an assumption that the embedding is optimal for
unknown-attack detection. The later evaluation is necessary because the
MLP was not trained with an explicit objective to separate every unseen
attack from benign activity. D. Isolation Forest Fitting The IF is
fitted to the hidden representations of the 433,538 benign training
records. Its configuration uses 300 estimators, a maximum sample size of
2,048 per tree, automatic con- tamination handling, random state 42, and
parallel execution. The detector is fitted without attack-labelled
records. In scikit- learn's decision-function convention, lower values
indicate greater abnormality; the project negates the returned decision
function so that a higher stored score corresponds to a more anomalous
observation. We denote this oriented score by s(z). The maximum sample
size is substantially smaller than the benign training set. This is
consistent with the subsampling design of Isolation Forest, but it means
that each tree is built from a subset of the training representations.
The ensemble combines the resulting isolation behaviour across trees.
The study does not tune the tree count or sample size through a nested
validation procedure; the parameters are fixed at the values above. E.
Calibration Rule and Operating Point Let SB cal be the anomaly scores
obtained from benign calibration records. The primary threshold is the
empirical 99th percentile: ?99 = Q0.99(SB cal) = 0.1633390. (6) A record
is flagged when s(z) \textgreater{} ?99. The percentile is selected from
a benign reference set, rather than from the distribution of
unknown-test attack scores. The intended interpretation is an initial
operating point that tolerates a small upper tail of benign calibration
scores. The 99th percentile does not guarantee a 1\% false-positive rate
on a later capture: shifts in vehicle state, collection conditions, or
the capture distribution can change the tail behaviour. A
97.5th-percentile threshold, ?97.5 = 0.0243811, is also reported. The
numerical threshold is calculated from benign calibration scores.
However, the decision to examine and emphasise this more permissive
quantile followed earlier in- spection of unknown-test metrics.
Therefore, the associated evaluation is exploratory. It is included to
show threshold sen- sitivity and to preserve the experimental record,
not to replace the primary result or to claim an unbiased improvement.
F. Evaluation Measures For the binary task, true positives (TP) are
attack records flagged by the detector; false positives (FP) are benign
records flagged as anomalous; true negatives (TN) are benign records
left unflagged; and false negatives (FN) are attack records left
unflagged. The metrics are Precision = TP TP + FP , (7) Recall = TP TP +
FN , (8) F1 = 2 Precision . Recall Precision + Recall , (9) FPR = FP FP
+ TN , (10) Accuracy = TP + TN TP + TN + FP + FN . (11) ROC-AUC and
average precision (reported as PR-AUC in the project outputs) are
calculated from continuous scores rather than binary predictions. This
distinction matters because threshold changes alter precision, recall,
F1, and FPR but do not alter the ranking of a fixed set of continuous
scores. In an imbalanced evaluation, a large true-negative count can
keep accuracy high even when many attacks are missed. Precision-- recall
measures are therefore reported with ROC-AUC and the confusion matrix
{[}9{]}. VI. EXPERIMENTAL SETUP AND BASELINES A. Implementation
Environment The implementation is written in Python and uses Py- Torch
for the MLP, scikit-learn for feature standardisation and Isolation
Forest, and pandas/NumPy for data handling. The recorded neural-network
run used CPU execution, batch size 4096, 25 epochs, and learning rate
0.001. The IF used 300 trees and 2,048 samples per tree. A fixed random
seed of 42 was specified in the principal scripts. The package list is
supplied separately in requirements.txt in the project repository. The
raw dataset is not redistributed with the research project. Reproduction
requires obtaining the dataset from its maintainers, arranging the
relevant capture files in the expected directory, and running the
preprocessing script to regenerate the processed CSV files. Trained
model checkpoints and local intermediate files are not part of the
source package. Accordingly, the manuscript reports the outputs obtained
in the project environment; it does not claim that a clean machine can
reproduce every numeric result without first obtaining the data and
NOVEMBER 2026 7 B. Original-Feature Isolation Forest The first
comparison model is an IF fitted directly to the original 13-dimensional
feature vectors. It uses the benign training rows for fitting and benign
calibration scores for the 99th-percentile threshold. The IF parameters
are 300 estima- tors, maximum sample size 2,048, automatic
contamination, and random state 42. This comparison helps answer whether
the learned representation changes anomaly-score ranking rel- ative to
the raw numeric inputs. It is not a comparison against the full set of
state-of-the-art CAN IDS methods. C. Energy-Based OOD Experiments The
project also executed a six-class MLP energy-score experiment. For
logits l(x), the score follows E(x) = -T log 6 ? k=1 exp(lk(x)/T), (12)
with temperature T = 1. The project evaluated a benign-calibration
threshold policy and a maximum-training- malicious-energy adaptation
inspired by the unknown- detection discussion in FedGAN-ID. Neither of
these executed energy experiments flagged any unknown attack records at
the reported operating points. This result is useful as a negative
comparison, but it does not establish that energy-based OOD methods
generally fail. It only describes these particular model outputs and
threshold policies on this evaluation file. D. Comparison Protocol All
methods in Table VI were evaluated on the same processed unknown-test
file. The primary comparison uses the 99th percentile for the
raw-feature IF and embedding IF. The 97.5th-percentile embedding result
is separated with an asterisk because it is exploratory. The energy
methods are included as executed project experiments. Direct comparison
to FedGAN-ID is intentionally excluded: FedGAN-ID uses a different
private vehicle dataset, different attack taxonomy, graph construction,
GAN, federated aggregation and privacy mechanisms. Without a common
benchmark and protocol, a numeric ranking between the two systems would
be mislead- ing. VII. RESULTS A. Primary Result at the 99th Percentile
At ?99 = 0.1633390, the embedding detector produces 3,956 true
positives, 8,937 false positives, 852,507 true neg- atives, and 16,869
false negatives. These counts give 30.68\% precision, 19.00\% recall,
and an F1-score of 23.47\%. The false-positive rate is 1.04\%, while
accuracy is 97.08\%. The accuracy value is not sufficient to describe
the detector: the number of benign records is much larger than the
number of attacks, and the detector misses more than four-fifths of the
attack records at this operating point. The relatively high ROC-AUC and
lower threshold-specific recall are not contradictory. ROC-AUC measures
the ranking of positive and negative records across possible cut-offs.
The TABLE III PRIMARY BINARY DETECTION RESULT AT THE 99TH-PERCENTILE
THRESHOLD Measure Value Threshold 0.1633390 Accuracy 97.08\% Precision
30.68\% Recall 19.00\% F1-score 23.47\% False-positive rate 1.04\%
ROC-AUC 0.9423 PR-AUC / average precision 0.4299 TP / FP 3,956 / 8,937
TN / FN 852,507 / 16,869 TABLE IV UNKNOWN ATTACK-FAMILY RESULTS AT THE
PRIMARY THRESHOLD Family Detected / total Recall DOUBLE 0 / 1,292 0.00\%
FUZZING 3,900 / 5,975 65.27\% INTERVAL 0 / 4,418 0.00\% SPEED 0 / 5,321
0.00\% SYSTEMATIC 56 / 1,043 5.37\% TRIPLE 0 / 2,776 0.00\%
99th-percentile rule chooses one cut-off from benign calibra- tion data.
If attack scores are often higher than benign scores but still fall
below that cut-off, the model can obtain a strong ranking measure while
returning a low recall at the selected operating point. The converse
operational issue is also visible: moving the cut-off lower identifies
more attack records but increases the number of benign records raised as
alerts. B. Per-Family Detection at the Primary Threshold The aggregate
recall of 19.00\% conceals large differences among attack families.
FUZZING contributes 3,900 detected records out of 5,975, giving 65.27\%
recall. SYSTEMATIC contributes 56 detections out of 1,043, or 5.37\%. No
records from DOUBLE, INTERVAL, SPEED, or TRIPLE exceed the primary
threshold. Table IV gives the counts and recalls. The zero-recall
outcomes should not be interpreted as proof that the corresponding
families are indistinguishable from benign traffic under every feature
representation. They indicate that, with this feature set, this trained
network, this IF configuration, and the primary threshold, none of the
evaluated records in those four families crossed the decision boundary.
The result identifies a concrete weakness in this configuration and
motivates family-specific analysis rather than a broad claim that unseen
CAN attacks are detected reliably. C. Threshold Sensitivity: Exploratory
97.5th Percentile The exploratory threshold is ?97.5 = 0.0243811. At
this value the detector flags 13,022 attacks and 18,905 benign records,
yielding 62.53\% recall, 40.79\% precision, 49.37\% F1, and a 2.19\%
false-positive rate. Recall increases by 43.53 percentage points
relative to the primary operating point; FPR increases by approximately
MANUSCRIPT, NOVEMBER 2026 8 Fig. 2. Recall by unknown attack family at
the primary 99th-percentile and exploratory 97.5th-percentile
thresholds. TABLE V PRIMARY AND EXPLORATORY OPERATING POINTS Metric 99th
(primary) 97.5th (exploratory) Threshold 0.1633390 0.0243811 Accuracy
97.08\% 96.97\% Precision 30.68\% 40.79\% Recall 19.00\% 62.53\%
F1-score 23.47\% 49.37\% FPR 1.04\% 2.19\% ROC-AUC 0.9423 0.9423 PR-AUC
0.4299 0.4299 TP 3,956 13,022 FP 8,937 18,905 FN 16,869 7,803 TN 852,507
842,539 number of detected attack records rises by 9,066, while the
number of false alerts rises by 9,968. This is a meaningful change in
the alert burden, not a free increase in detector quality. The
97.5th-percentile result may be attractive where missed attacks are
considered more costly than additional alerts, but that operational
preference would need to be defined before examining the test labels. In
this study, the unknown-test file was inspected during earlier threshold
exploration. As a result, the selected 97.5th- percentile point is
reported transparently as exploratory and should not be used as the
headline validated performance figure. At this exploratory operating
point, family-wise recall is 9.06\% for DOUBLE, 80.52\% for FUZZING,
96.42\% for INTERVAL, 60.27\% for SPEED, 34.32\% for SYSTEMATIC, and
9.69\% for TRIPLE. The gain is therefore not evenly distributed.
INTERVAL and SPEED change from zero recall at the primary threshold to
high or moderate recall at the exploratory threshold, while DOUBLE and
TRIPLE remain difficult. Fig. 2 presents both sets of family-wise
results. The ROC-AUC and PR-AUC values remain the same be- tween the two
threshold rows because both are computed from the same continuous score
vector on the same evaluation records. Threshold choice changes the
binary predictions, not the score ranking. Reporting the same AUC values
beside different F1 values is therefore expected; it is not a separate
improvement in ranking. Fig. 3. Confusion matrices for the primary and
exploratory operating points. The exploratory panel is included to
illustrate threshold sensitivity and is not an independent validation
result. Fig. 4. Threshold-specific F1, recall, and FPR for the two
Isolation Forest configurations. The 97.5th-percentile point is
exploratory. D. Comparison with the Original-Feature Baseline The IF
trained on the original 13 features obtains F1 28.94\%, recall 24.47\%,
and FPR 1.08\% at its 99th-percentile threshold. The embedding-based IF
obtains F1 23.47\%, recall 19.00\%, and FPR 1.04\%. Thus, by the
selected threshold- specific metrics, the original-feature IF performs
better on F1 and recall, while the embedding detector has a slightly
lower FPR. The ranking measures show a different pattern. The original-
feature IF has ROC-AUC 0.6303 and PR-AUC 0.2006, whereas the embedding
IF has ROC-AUC 0.9423 and PR- AUC 0.4299. The embedding provides a much
stronger or- dering of the evaluated benign and attack records according
to these area measures. This does not make the embedding detector
unambiguously superior. It means its score distribu- tion separates many
positive and negative examples in ranking terms, but the selected
99th-percentile cut-off is conservative for this test distribution. A
detector selection decision should take account of both ranking quality
and the alert rate at the intended operating point. The energy-based
experiments did not detect unknown records at the chosen operating
points. Their near-zero FPR is not a useful result by itself because it
is accompanied by zero recall. In practical intrusion detection, a
detector that never raises an attack alert is not effective merely
because it rarely alarms on benign traffic. This negative result is
included because it influenced the model-selection process; it should
not be generalised to all energy-based OOD methods. Page 10 of 14 - AI
PROJECT MANUSCRIPT, NOVEMBER 2026 9 TABLE VI EXECUTED METHOD COMPARISONS
ON THE SAME UNKNOWN-TEST FILE Method F1 Recall FPR ROC-AUC Raw-feature
IF, 99th 28.94\% 24.47\% 1.08\% 0.6303 Embedding IF, 99th 23.47\%
19.00\% 1.04\% 0.9423 Embedding IF, 97.5th* 49.37\% 62.53\% 2.19\%
0.9423 Energy OOD 0.00\% 0.00\% 0.001\% 0.2325 Max-training-energy
0.00\% 0.00\% 0.0017\% 0.2325 * Exploratory threshold: the evaluation
file had been inspected during earlier threshold analysis. The energy
rows describe only the two executed project configurations. Fig. 5.
Median anomaly score by ground-truth group. Higher scores are more
anomalous under the stored score convention. E. Anomaly-Score
Distribution by Family The median anomaly score differs across the
evaluated la- bels. The benign median is -0.131440. The medians for DOU-
BLE and TRIPLE are -0.108853 and 0.006674, respectively, while FUZZING
has a much higher median of 0.297320. INTERVAL and SPEED have medians of
0.044970 and 0.033780, and SYSTEMATIC has a median of -0.029677. These
summaries are consistent with the observed recall pat- tern: FUZZING
scores tend to lie further into the anomalous tail, whereas the score
distributions for DOUBLE and TRIPLE overlap more strongly with the
benign region. A median is only one point in a distribution. It does not
describe the tails, spread, or proportion of records above a selected
threshold. Therefore, the median plot is used as a descriptive
diagnostic, not as a substitute for the family-wise confusion counts. In
particular, a family may have a median below the threshold but still
contain a high-scoring subset, and a median above the threshold does not
imply that every record is flagged. F. Timing and Payload Diagnostics
The processed data also provide timing and payload sum- maries that help
interpret the family-specific outcomes. The benign median inter-arrival
time is 0.0138 ms, compared with 0.8030 ms for DOUBLE and 0.7639 ms for
TRIPLE. Median payload length is 8 bytes for benign records, 7 for
DOUBLE, and 5 for TRIPLE. FUZZING has a median inter-arrival time of
0.8988 ms and a very large median same-identifier inter- arrival time in
the processed summary. INTERVAL shows a median inter-arrival time of
0.3589 ms and a same-identifier inter-arrival time of 0.3641 ms. SPEED
records show median inter-arrival and same-identifier inter-arrival
times of 0.8707 ms and 2.9240 ms, respectively. These are marginal
summaries. They do not demonstrate that a specific timing feature caused
a detection or miss. The MLP may combine all input dimensions
nonlinearly, and the IF acts on the 64-dimensional representation rather
than on an individual feature. A controlled ablation study would be
needed to attribute the difference in family recall to timing, payload,
or identifier information. No such ablation is claimed in this
manuscript. VIII. DISCUSSION A. Ranking Performance and the Selected
Threshold The embedding-based detector obtains ROC-AUC 0.9423 and PR-AUC
0.4299. At the primary threshold, however, recall is 19.00\%. These
values describe different properties of the detector. The AUC measures
how the scores order benign and attack records over a range of cut-offs;
recall measures the fraction of attacks detected at the single threshold
used here. A strong ranking result therefore does not guarantee useful
decisions at a particular operating point {[}9{]}. The exploratory
97.5th-percentile cut-off raises recall from 19.00\% to 62.53\%. The
false-positive rate also rises, from 1.04\% to 2.19\%. On this test
file, that change corresponds to 9,066 additional detected attacks and
9,968 additional false alarms. The 97.5th-percentile result is not an
independent validation result: the test file had already been inspected
during threshold exploration. A deployment threshold should instead be
selected using a separate validation capture and a predeclared
false-alarm budget, then held fixed for the final evaluation. B.
Detection Differences Across Attack Families The primary threshold
detects 65.27\% of FUZZING records and 5.37\% of SYSTEMATIC records.
None of the evaluated DOUBLE, INTERVAL, SPEED, or TRIPLE records crosses
this threshold. At the exploratory threshold, INTERVAL and SPEED recall
increase substantially, while DOUBLE and TRIPLE remain below 10\%. The
score summaries show that the family distributions are not identical.
For example, FUZZING has a higher median anomaly score than benign
traffic, whereas the DOUBLE median is close to the benign median. This
observation is consistent with the measured recall differences, but it
does not identify the feature responsible for them. The MLP uses all 13
input features, and the IF operates on the resulting 64-dimensional
representation. Feature ablations or sequence- level experiments would
MANUSCRIPT, NOVEMBER 2026 10 payload, or identifier information explains
the family-wise outcomes. C. Raw Features and Learned Representations
The original-feature IF obtains an F1-score of 28.94\% and recall of
24.47\% at its 99th-percentile threshold. The embed- ding IF obtains
23.47\% F1 and 19.00\% recall. In contrast, the embedding IF has higher
ROC-AUC (0.9423 versus 0.6303) and PR-AUC (0.4299 versus 0.2006). Thus,
the embedding changes the score ordering, but this does not translate
into better binary decisions at the selected threshold. The results do
not support describing the embedding detec- tor as uniformly better. Its
ranking metrics are higher, while the raw-feature baseline has better F1
and recall at the primary operating point. Which configuration is
preferable depends on the intended operating requirement and the
threshold policy. Both ranking metrics and threshold-specific measures
should therefore be reported. D. Energy-Based OOD Experiments The two
executed energy-score experiments did not flag any unknown attack
records at their reported thresholds. This is the outcome for the
particular MLP, score definition, calibration data, and thresholds
tested here. It does not show that energy- based OOD detection is
generally ineffective. One possible explanation is that the MLP assigns
unknown records logits similar to those of known classes, resulting in
energy scores that do not exceed the selected cut-off. The current
experiments do not establish this explanation: they do not include a
family-wise analysis of logit confidence or a controlled comparison of
temperature and calibration choices. The implementation is also not the
full FedGAN-ID system, which combines energy-based OOD detection with
graph construction, GAN-based modelling, and federated learning {[}5{]}.
No conclusion about the performance of that complete system can be drawn
from these local experiments. E. Calibration Under Distribution Shift
The 99th-percentile calibration rule is simple and repro- ducible: it
fixes the threshold from the upper tail of benign calibration scores.
Its weakness is that the calibration and test score distributions may
differ. A new driving session, a different vehicle, changed bus load,
maintenance activity, or a modified capture setup may shift the benign
score distribution. A threshold that corresponds to one percent of
benign calibra- tion scores can therefore produce a different
false-positive rate elsewhere. This study observes a 1.04\% FPR on the
processed test file, which is close to but not guaranteed by the 99th-
percentile calibration rule. An operational calibration procedure should
test stability over multiple benign sessions and vehicles. It should
also include a method for handling drift. Updating a threshold
automatically is not straightforward: if the data used for recalibration
contain undetected attacks, the benign reference region may be
contaminated. Future work should investigate drift detection and
controlled recalibration with trusted benign intervals rather than
treating a single calibration percentile as a universal constant. F.
Implications for a Vehicle IDS The implemented pipeline can be
understood as a research prototype for record-level anomaly scoring. It
has not been evaluated on an automotive electronic control unit, a
produc- tion gateway, or a hardware-in-the-loop setup. The reported
training time and offline processing do not establish real- time
inference feasibility. A deployment study would need to measure
per-frame latency, throughput under bus load, memory consumption, model
update cost, and the effect of logging and preprocessing on an embedded
processor. It would also need to test whether timing features can be
calculated with consistent semantics in the target gateway. A further
distinction concerns alerts. The current output is binary and does not
identify the originating ECU or ex- plain which byte or timing feature
caused the score. For use by a vehicle-security engineer, an alert would
ideally in- clude contextual evidence: identifier, capture time, local
score, threshold, surrounding frames, and a confidence or severity
interpretation. That extension must be designed carefully so that
explanation does not create false confidence or disclose sensitive
vehicle information. IX. VALIDITY CONSIDERATIONS AND LIMITATIONS Several
limitations determine how far the results can be generalised. A.
Evaluation-Set Reuse The unknown-test file was examined during earlier
ex- ploratory threshold analysis. Even though the threshold value is
calculated from benign calibration scores, the decision to emphasise the
97.5th percentile followed inspection of test performance. This creates
a model-selection concern: the exploratory result is no longer an
untouched estimate of generalisation. The 99th-percentile operating
point is retained as the primary result because it was the initially
selected operating point. A new capture that has not been consulted
during development is required for an independent estimate. B. Class
Prevalence and Sampling Benign traffic was subsampled while attack
records were retained. The unknown-test file has an attack prevalence of
approximately 2.36\%. Precision and accuracy depend on class prevalence,
so they should not be transferred directly to a deployment environment
with a different benign-to-attack ratio. Recall and FPR are also subject
to dataset and capture shift. The sampling policy must be disclosed when
using the reported precision and accuracy in a comparison. C. Limited
Dataset and Capture Scope The experiments use one configured known-class
train- ing/calibration workflow and one designated unknown-attack test
file. They do not establish performance over all vehicles, software
versions, road conditions, and bus configurations. The
can-train-and-test paper was designed to support broader comparisons
across attack types and vehicles, but this project has not run all of
its subsets {[}8{]}. Consequently, no universal cross-vehicle claim is
NOVEMBER 2026 11 D. Single Configuration and Randomness The reported
model uses a fixed seed and one configured training run. There are no
multi-seed means, standard devi- ations, confidence intervals, or
statistical significance tests. The IF itself uses random partitioning.
Although fixed random states make the recorded run more repeatable, they
do not show how stable the performance is across different seeds. A
stronger study would repeat the full pipeline using several seeds and
report variation in both aggregate and family-wise metrics. E.
Representation and Feature Limitations The MLP is trained for six-class
classification, not for open- set representation learning. The hidden
space may suppress variations that do not help separate the known
classes. The 13-dimensional input contains only a small set of per-frame
and local timing attributes. It does not include a learned graph of
message transitions, a longer sequence window, physical- layer
information, or explicit transmitter identity. Numeric CAN identifier
encoding may also impose distances that do not reflect semantic
relations among message IDs. F. No Complete FedGAN-ID Reproduction The
work does not implement the GAN, federated aggre- gation, differential
privacy, message-graph construction, MTF, READ, or the second-level
classifier described in FedGAN- ID. It also does not evaluate privacy
leakage, communication overhead, client heterogeneity, or attack
detection latency on a real vehicle. The use of FedGAN-ID as motivation
must not be interpreted as proof that the present project reproduces the
article's architecture or validates its reported results {[}5{]}. X.
REPRODUCIBILITY AND REPORTING CHECKLIST To make a future rerun easier to
audit, the following items should accompany the paper: (i) the exact
repository commit used for the submitted manuscript; (ii) the dataset
version and capture filenames; (iii) the preprocessing script and its
sampling settings; (iv) the Python and library versions; (v) the MLP
checkpoint and scaler or an instruction to regenerate them; (vi) the
seed and training parameters; (vii) the exact calibration file and
benign-row filtering rule; (viii) the IF configuration and threshold
percentile; and (ix) the scripts that generate every reported table and
figure. The raw dataset should be obtained from its maintainers rather
than redis- tributed without checking its licence. A reproducibility
package should distinguish source code from generated outputs. The
scripts and configuration are necessary to reproduce the analysis, while
the CSV result summaries and figures are convenient evidence of the par-
ticular run. A clean installation test should start from the repository
root, recreate the environment, run preprocessing, train the model,
extract embeddings, fit the detector, and execute evaluation. A syntax
check alone is not a runtime reproduction. The current manuscript
therefore describes the implementation and its recorded outputs without
claiming that a third party has independently rerun the complete
workflow. XI. FUTURE WORK A first priority is to obtain an untouched
evaluation capture. The threshold percentile, feature set, and model
configuration should be frozen using training and validation data only.
The final capture should be opened only after the complete decision rule
is fixed. This would provide a cleaner estimate of whether the
99th-percentile operating point transfers to a new session and whether a
more permissive threshold is justified by a predeclared alert budget. A
second direction is to evaluate the effect of the input representation.
Ablations could remove payload bytes, CAN identifier, or each timing
feature group in turn. A temporal model could process a sequence of
consecutive records rather than independent row vectors. A graph-based
model could represent message transitions or co-occurrence patterns.
These experiments would help test the hypothesis that DOUBLE and TRIPLE
are difficult because their local feature distributions overlap with
benign traffic. Such an explanation should only be accepted if the
controlled experiments support it. Third, the threshold should be
evaluated as an operational parameter. Instead of selecting a percentile
after observing test results, the study could choose a target benign
alert rate on a validation set, freeze the threshold, and report its
performance across separate captures. Calibration stability should be
mea- sured by vehicle, session, and operating condition. The model
should also be assessed using precision--recall curves, family- wise
metrics, and confidence intervals across repeated training runs. Fourth,
the energy-based baseline should be investigated with a controlled study
of logit distributions and threshold policies. The present zero-recall
outcome suggests that the current energy score is not separating the
unknown attack examples at the tested threshold. It does not reveal
whether the issue is the MLP, the score definition, the temperature, the
calibration population, or the attack distribution. These factors should
be varied one at a time on validation data before any claim is made.
Finally, federated and privacy-preserving extensions could be explored
if multiple vehicle datasets become available. Such an extension would
require a true multi-client setup, explicit local training and
aggregation, a documented privacy mechanism and budget, and analysis of
the effect of privacy noise on anomaly detection. Those are substantial
additional experiments; they are not minor modifications to the current
IF pipeline. XII. CONCLUSION The MLP--IF pipeline was evaluated on one
designated unknown-attack capture from CAN-train-and-test v1.5. At the
primary 99th-percentile threshold, it detected 3,956 of 20,825 attack
records. This gives 19.00\% recall, 30.68\% precision, and 23.47\% F1,
with a 1.04\% false-positive rate. The ROC- AUC was 0.9423 and the
PR-AUC was 0.4299. Detection was concentrated in FUZZING (65.27\%
recall); no DOUBLE, INTERVAL, SPEED, or TRIPLE records crossed the
NOVEMBER 2026 12 The comparison with the raw-feature IF produced a mixed
result. The embedding detector ranked the evaluated records better
according to ROC-AUC and PR-AUC, but the raw- feature detector achieved
higher F1 and recall at the selected operating point. The exploratory
97.5th-percentile threshold increased recall to 62.53\% and F1 to
49.37\%, at a 2.19\% false- positive rate. Because the unknown-test file
was inspected during threshold exploration, this result should not be
treated as an independent estimate of generalisation. The experiments
describe this MLP--IF implementation only. They do not reproduce the
GAN, federated aggregation, differential privacy, graph construction, or
other components of FedGAN-ID. A stronger evaluation requires a new
capture that has not been used during development, repeated training
runs, and a threshold fixed using validation data before the final test.
Until then, the reported measurements apply to the dataset split and
configuration described in this paper. ACKNOWLEDGMENT The authors thank
Dr.~Bhale Pradeepkumar Gajendra for guidance during the research
project. The authors are respon- sible for verifying the correctness of
the experiments, interpret- ing the results, and ensuring that the final
manuscript complies with the institution's research and authorship
policies. REFERENCES {[}1{]} Robert Bosch GmbH, CAN Specification
Version 2.0, Robert Bosch GmbH, 1991, technical specification. {[}2{]}
S. Rajapaksha, H. Kalutarage, M. O. Al-Kadri, A. Petrovski, G. Madzudzo,
and M. Cheah, ``AI-based intrusion detection systems for in-vehicle
networks: A survey,'' ACM Computing Surveys, vol.~55, no. 11, pp.~1--40,
2023. {[}3{]} K.-T. Cho and K. G. Shin, ``Fingerprinting electronic
control units for vehicle intrusion detection,'' in Proceedings of the
25th USENIX Security Symposium. USENIX Association, 2016, pp.~911--927.
{[}Online{]}. Available: https://www.usenix.org/conference/
usenixsecurity16/technical-sessions/presentation/cho {[}4{]} M. Muter
and N. Asaj, ``Entropy-based anomaly detection for in-vehicle
networks,'' in 2011 IEEE Intelligent Vehicles Symposium, 2011, pp.
1110--1115. {[}5{]} M. Chen, B. Wu, H. Sun, and Z. Wang, ``FedGAN-ID:
Federated- learning-based intrusion detection for in-vehicle network
using GANs,'' IEEE Internet of Things Journal, vol.~12, no. 17, pp.~36
155--36 167, Sep.~2025. {[}6{]} M. Althunayyan, A. Javed, and O. Rana,
``A survey of learning- based intrusion detection systems for in-vehicle
networks,'' Computer Networks, vol.~277, p.~112031, 2026. {[}7{]} F. T.
Liu, K. M. Ting, and Z.-H. Zhou, ``Isolation forest,'' in Proceedings of
the 8th IEEE International Conference on Data Mining, 2008, pp.
413--422. {[}8{]} B. Lampe and W. Meng, ``can-train-and-test: A curated
CAN dataset for automotive intrusion detection,'' Computers and
Security, vol.~140, p.~103777, 2024. {[}9{]} J. Davis and M. Goadrich,
``The relationship between precision-recall and ROC curves,'' in
Proceedings of the 23rd International Conference on Machine Learning.
ACM, 2006, pp.~233--240. {[}10{]} B. Lampe and W. Meng,
``CAN-train-and-test-v1.5 dataset repository,''
https://bitbucket.org/brooke-lampe/can-train-and-test-v1.5/src/master/,

\end{document}
